% Options for packages loaded elsewhere
\PassOptionsToPackage{unicode,hidelinks}{hyperref}
\PassOptionsToPackage{hyphens}{url}
%
\documentclass[
  letterpaper,
  a4paper,
  12pt,
  titlepage,
  oneside,
  openany]{book}
\usepackage{amsmath,amssymb}
\usepackage{lmodern}
\usepackage{iftex}
\ifPDFTeX
  \usepackage[T1]{fontenc}
  \usepackage[utf8]{inputenc}
  \usepackage{textcomp} % provide euro and other symbols
\else % if luatex or xetex
  \usepackage{unicode-math}
  \defaultfontfeatures{Scale=MatchLowercase}
  \defaultfontfeatures[\rmfamily]{Ligatures=TeX,Scale=1}
\fi

% Use upquote if available, for straight quotes in verbatim environments
\IfFileExists{upquote.sty}{\usepackage{upquote}}{}
\IfFileExists{microtype.sty}{% use microtype if available
  \usepackage[]{microtype}
  \UseMicrotypeSet[protrusion]{basicmath} % disable protrusion for tt fonts
}{}
\setlength{\parskip}{0pt}
\usepackage{xcolor}

\usepackage[paperheight = 29.7cm,paperwidth = 21cm,outer = 1.5cm,inner =
2.5cm,top = 2cm,bottom = 2cm]{geometry}

\setlength{\emergencystretch}{3em} % prevent overfull lines
\setcounter{secnumdepth}{5}


\pagestyle{plain}

%%% This section was taken from `https://aty.sdsu.edu/bibliog/latex/floats.html`
% Alter some LaTeX defaults for better treatment of figures:
% See p.105 of "TeX Unbound" for suggested values.
% See pp. 199-200 of Lamport's "LaTeX" book for details.
%   General parameters, for ALL pages:
\renewcommand{\topfraction}{0.9}	% max fraction of floats at top
\renewcommand{\bottomfraction}{0.8}	% max fraction of floats at bottom
%   Parameters for TEXT pages (not float pages):
\setcounter{topnumber}{2}
\setcounter{bottomnumber}{2}
\setcounter{totalnumber}{4}     % 2 may work better
\setcounter{dbltopnumber}{2}    % for 2-column pages
\renewcommand{\dbltopfraction}{0.9}	% fit big float above 2-col. text
\renewcommand{\textfraction}{0.07}	% allow minimal text w. figs
%   Parameters for FLOAT pages (not text pages):
\renewcommand{\floatpagefraction}{0.7}	% require fuller float pages
% N.B.: floatpagefraction MUST be less than topfraction !!
\renewcommand{\dblfloatpagefraction}{0.7}	% require fuller float pages


\usepackage{color}
\usepackage{fancyvrb}
\newcommand{\VerbBar}{|}
\newcommand{\VERB}{\Verb[commandchars=\\\{\}]}
\DefineVerbatimEnvironment{Highlighting}{Verbatim}{commandchars=\\\{\}}
% Add ',fontsize=\small' for more characters per line
\usepackage{framed}
\definecolor{shadecolor}{RGB}{255,255,255}
\newenvironment{Shaded}{\begin{snugshade}}{\end{snugshade}}
\newcommand{\AlertTok}[1]{\textcolor[rgb]{0.00,0.00,0.00}{#1}}
\newcommand{\AnnotationTok}[1]{\textcolor[rgb]{0.00,0.00,0.00}{\textit{#1}}}
\newcommand{\AttributeTok}[1]{\textcolor[rgb]{0.00,0.00,0.00}{#1}}
\newcommand{\BaseNTok}[1]{\textcolor[rgb]{0.00,0.00,0.00}{#1}}
\newcommand{\BuiltInTok}[1]{\textcolor[rgb]{0.00,0.00,0.00}{\textbf{#1}}}
\newcommand{\CharTok}[1]{\textcolor[rgb]{0.00,0.00,0.00}{#1}}
\newcommand{\CommentTok}[1]{\textcolor[rgb]{0.00,0.00,0.00}{#1}}
\newcommand{\CommentVarTok}[1]{\textcolor[rgb]{0.00,0.00,0.00}{#1}}
\newcommand{\ConstantTok}[1]{\textcolor[rgb]{0.00,0.00,0.00}{#1}}
\newcommand{\ControlFlowTok}[1]{\textcolor[rgb]{0.00,0.00,0.00}{\textbf{#1}}}
\newcommand{\DataTypeTok}[1]{\textcolor[rgb]{0.00,0.00,0.00}{\textbf{#1}}}
\newcommand{\DecValTok}[1]{\textcolor[rgb]{0.00,0.00,0.00}{#1}}
\newcommand{\DocumentationTok}[1]{\textcolor[rgb]{0.00,0.00,0.00}{\textit{#1}}}
\newcommand{\ErrorTok}[1]{\textcolor[rgb]{0.00,0.00,0.00}{#1}}
\newcommand{\ExtensionTok}[1]{\textcolor[rgb]{0.00,0.00,0.00}{#1}}
\newcommand{\FloatTok}[1]{\textcolor[rgb]{0.00,0.00,0.00}{\textbf{#1}}}
\newcommand{\FunctionTok}[1]{\textcolor[rgb]{0.00,0.00,0.00}{#1}}
\newcommand{\ImportTok}[1]{\textcolor[rgb]{0.00,0.00,0.00}{\textbf{#1}}}
\newcommand{\InformationTok}[1]{\textcolor[rgb]{0.00,0.00,0.00}{\textit{#1}}}
\newcommand{\KeywordTok}[1]{\textcolor[rgb]{0.00,0.00,0.00}{\textbf{#1}}}
\newcommand{\NormalTok}[1]{\textcolor[rgb]{0.00,0.00,0.00}{#1}}
\newcommand{\OperatorTok}[1]{\textcolor[rgb]{0.00,0.00,0.00}{#1}}
\newcommand{\OtherTok}[1]{\textcolor[rgb]{0.00,0.00,0.00}{#1}}
\newcommand{\PreprocessorTok}[1]{\textcolor[rgb]{0.00,0.00,0.00}{#1}}
\newcommand{\RegionMarkerTok}[1]{\textcolor[rgb]{0.00,0.00,0.00}{#1}}
\newcommand{\SpecialCharTok}[1]{\textcolor[rgb]{0.00,0.00,0.00}{#1}}
\newcommand{\SpecialStringTok}[1]{\textcolor[rgb]{0.00,0.00,0.00}{#1}}
\newcommand{\StringTok}[1]{\textcolor[rgb]{0.00,0.00,0.00}{#1}}
\newcommand{\VariableTok}[1]{\textcolor[rgb]{0.00,0.00,0.00}{\textbf{#1}}}
\newcommand{\VerbatimStringTok}[1]{\textcolor[rgb]{0.00,0.00,0.00}{#1}}
\newcommand{\WarningTok}[1]{\textcolor[rgb]{0.00,0.00,0.00}{\textit{#1}}}

\providecommand{\tightlist}{%
  \setlength{\itemsep}{0pt}\setlength{\parskip}{0pt}}\usepackage{longtable,booktabs,array}
\usepackage{calc} % for calculating minipage widths
% Correct order of tables after \paragraph or \subparagraph
\usepackage{etoolbox}
\makeatletter
\patchcmd\longtable{\par}{\if@noskipsec\mbox{}\fi\par}{}{}
\makeatother
% Allow footnotes in longtable head/foot
\IfFileExists{footnotehyper.sty}{\usepackage{footnotehyper}}{\usepackage{footnote}}
\makesavenoteenv{longtable}
\usepackage{graphicx}
\makeatletter
\def\maxwidth{\ifdim\Gin@nat@width>\linewidth\linewidth\else\Gin@nat@width\fi}
\def\maxheight{\ifdim\Gin@nat@height>\textheight\textheight\else\Gin@nat@height\fi}
\makeatother
% Scale images if necessary, so that they will not overflow the page
% margins by default, and it is still possible to overwrite the defaults
% using explicit options in \includegraphics[width, height, ...]{}
\setkeys{Gin}{width=\maxwidth,height=\maxheight,keepaspectratio}
% Set default figure placement to htbp
\makeatletter
\def\fps@figure{htbp}
\makeatother
\newlength{\cslhangindent}
\setlength{\cslhangindent}{1.5em}
\newlength{\csllabelwidth}
\setlength{\csllabelwidth}{3em}
\newlength{\cslentryspacingunit} % times entry-spacing
\setlength{\cslentryspacingunit}{\parskip}
\newenvironment{CSLReferences}[2] % #1 hanging-ident, #2 entry spacing
 {% don't indent paragraphs
  \setlength{\parindent}{0pt}
  % turn on hanging indent if param 1 is 1
  \ifodd #1
  \let\oldpar\par
  \def\par{\hangindent=\cslhangindent\oldpar}
  \fi
  % set entry spacing
  \setlength{\parskip}{#2\cslentryspacingunit}
 }%
 {}
\usepackage{calc}
\newcommand{\CSLBlock}[1]{#1\hfill\break}
\newcommand{\CSLLeftMargin}[1]{\parbox[t]{\csllabelwidth}{#1}}
\newcommand{\CSLRightInline}[1]{\parbox[t]{\linewidth - \csllabelwidth}{#1}\break}
\newcommand{\CSLIndent}[1]{\hspace{\cslhangindent}#1}

% Additional required packages
\usepackage{plain}
\usepackage{setspace}
\usepackage{titlesec}

% Remove empty vertical space around quotes
\renewcommand{\quote}{\list{}{\rightmargin=\leftmargin\topsep=0pt}\item\relax}

% Define a changemargin block
\def\changemargin#1#2{\list{}{\rightmargin#2\leftmargin#1}\item[]}
\let\endchangemargin=\endlist
% Change margin for the Shaded environment
\renewenvironment{Shaded}{
\begin{changemargin}{1cm}{1cm}
  \begin{snugshade}
}{
  \end{snugshade}
\end{changemargin}
}
% Change line number size
\renewcommand{\theFancyVerbLine}{\small\arabic{FancyVerbLine}}
\makeatletter
\makeatother
\makeatletter
\@ifpackageloaded{bookmark}{}{\usepackage{bookmark}}
\makeatother
\makeatletter
\@ifpackageloaded{caption}{}{\usepackage{caption}}
\AtBeginDocument{%
\ifdefined\contentsname
  \renewcommand*\contentsname{Table of contents}
\else
  \newcommand\contentsname{Table of contents}
\fi
\ifdefined\listfigurename
  \renewcommand*\listfigurename{List of Figures}
\else
  \newcommand\listfigurename{List of Figures}
\fi
\ifdefined\listtablename
  \renewcommand*\listtablename{List of Tables}
\else
  \newcommand\listtablename{List of Tables}
\fi
\ifdefined\figurename
  \renewcommand*\figurename{Figure}
\else
  \newcommand\figurename{Figure}
\fi
\ifdefined\tablename
  \renewcommand*\tablename{Table}
\else
  \newcommand\tablename{Table}
\fi
}
\@ifpackageloaded{float}{}{\usepackage{float}}
\floatstyle{ruled}
\@ifundefined{c@chapter}{\newfloat{codelisting}{h}{lop}}{\newfloat{codelisting}{h}{lop}[chapter]}
\floatname{codelisting}{Listing}
\newcommand*\listoflistings{\listof{codelisting}{List of Listings}}
\makeatother
\makeatletter
\@ifpackageloaded{caption}{}{\usepackage{caption}}
\@ifpackageloaded{subcaption}{}{\usepackage{subcaption}}
\makeatother
\makeatletter
\makeatother

\ifLuaTeX
\usepackage[bidi=basic]{babel}
\else
\usepackage[bidi=default]{babel}
\fi
\babelprovide[main,import]{english}
% get rid of language-specific shorthands (see #6817):
\let\LanguageShortHands\languageshorthands
\def\languageshorthands#1{}
\ifLuaTeX
  \usepackage{selnolig}  % disable illegal ligatures
\fi
\usepackage{csquotes}
\IfFileExists{bookmark.sty}{\usepackage{bookmark}}{\usepackage{hyperref}}
\IfFileExists{xurl.sty}{\usepackage{xurl}}{} % add URL line breaks if available
\urlstyle{same} % disable monospaced font for URLs
\hypersetup{
  pdftitle={OpenRFSense - Explorative data analysis using software-defined radios},
  pdfauthor={Leonardo Baldin},
  pdflang={en},
  hidelinks,
  pdfcreator={LaTeX via pandoc}}
\title{OpenRFSense - Explorative data analysis using software-defined
radios}
\author{Leonardo Baldin}
\date{2023-03-15}

\begin{document}
\frontmatter

\thispagestyle{empty}

\begin{titlepage}
  \begin{center}
    \begin{figure}[h!]
      \centerline{\includegraphics[width=0.6\textwidth]{\_extensions/unitn-thesis/assets/marchio.pdf}}
    \end{figure}

    \vspace{2 cm}

    \LARGE{Department of Information Engineering and Computer Science\\}

    \vspace{1 cm}
    \Large{Bachelor's Degree in\\
      Computer Science
    }

    \vspace{2 cm}
    \Large\textsc{Final Dissertation\\}
    \vspace{1 cm}
    \Huge\textsc{OpenRFSense - Explorative data analysis using
software-defined radios\\}
    

    \vspace{2 cm}
    \begin{tabular*}{\textwidth}{ c @{\extracolsep{\fill}} c }
    \Large{Supervisor} & \Large{Student}\\
    \Large{Fabrizio Granelli}& \Large{Leonardo Baldin}\\
    \end{tabular*}

    \vspace{2 cm}

    \Large{Academic year 2021/2022}
  \end{center}
\end{titlepage}

\addtocounter{page}{-1}

\mainmatter
\renewcommand*\contentsname{Table of contents}
{
\setcounter{tocdepth}{2}
\tableofcontents
}
\clearpage

\begingroup
  % Save the \chapter command before patching
  \let\OldChapter\chapter
  % Patch the \chapter command to not create a new page
  \makeatletter
  \patchcmd\chapter
    {\if@openright\cleardoublepage\else\clearpage\fi}
    {\par}
    {}{}
  \makeatother
  % Restore the command behavior when \appendix is called
  \preto\appendix
    {\renewcommand{\chapter}{\OldChapter}}
    {}{}

  \titleformat{\chapter}
    {\normalfont\Huge\bfseries}{\thechapter}{1em}{}

  \titlespacing*{\chapter}{0pt}{1.5cm}{0.05cm}
  \titlespacing*{\section}{0pt}{0.5cm}{0.05cm}
  \titlespacing*{\subsection}{0pt}{0.25cm}{0.05cm}
\bookmarksetup{startatroot}

\hypertarget{abstract}{%
\chapter*{Abstract}\label{abstract}}
\addcontentsline{toc}{chapter}{Abstract}

\markboth{Abstract}{Abstract}

Placeholder for the abstract.

\bookmarksetup{startatroot}

\hypertarget{introduction}{%
\chapter{Introduction}\label{introduction}}

\begin{figure}[h!]

{\centering \includegraphics[width=0.8\textwidth,height=\textheight]{assets/SchemaExtract2.gif}

}

\caption{SchemaExtract}

\end{figure}

\hypertarget{background}{%
\section{Background}\label{background}}

Truly understanding customers is key to a successful business. Companies
increasingly leverage data analytics to uncover insights into their
customers' needs, preferences, and socio-demographic profiles. This
information is essential for enhancing user experiences, differentiating
from competitors, and creating products that satisfy market demands.
Moreover, understanding customers improves strategic planning by
facilitating informed decisions about market positioning, pricing, and
promotional strategies. By aligning marketing strategies with consumer
needs, businesses can achieve more successful and impactful outcomes.

Traditional approaches to understanding customers, such as surveys,
focus groups, and market research reports, have long been relied upon.
While these methods can provide useful insights, they have significant
limitations. Surveys, for instance, often suffer from low response rates
and biased results. Focus groups may not accurately represent the entire
customer base, leading to incomplete or skewed findings. Market research
reports, while detailed, are frequently expensive and time-consuming to
produce. Additionally, these traditional methods typically provide only
a static view of customer preferences, failing to capture the dynamic
changes in consumer behavior over time.

To address these limitations, businesses are increasingly turning to
more innovative approaches. One such approach is the analysis of
customer-generated content, such as product reviews. Product reviews
offer valuable insights into customer preferences and sentiments.
Various techniques, including sentiment analysis, topic modeling, and
natural language processing, can be employed to extract meaningful
information from this user-generated content.

\hypertarget{research-objectives}{%
\section{Research Objectives}\label{research-objectives}}

Traditionally, product reviews have been analyzed primarily on feedback
about the product, such as customer opinions, satisfaction levels, and
sentiments on specific features. However, insights into customers'
personal attributes, behaviors, and usage contexts have not been
thoroughly explored. Addressing this gap presents a valuable opportunity
to leverage user-generated content for a deeper understanding of
customers.

The extraction of personal attributes from product reviews is the focus
of this research. Personal attributes are specific characteristics that
uniquely describe an individual's traits, lifestyle, and circumstances,
such as physical appearance, profession, household composition, and
living situation. These attributes could include details such as hair
type, skin condition, profession, number of family members or pets, and
whether they live in an apartment, townhouse or single-family home. By
extracting personal attributes from product reviews, we aim to gain a
deeper understanding of the characteristics of customers.

\hypertarget{related-work}{%
\section{Related work}\label{related-work}}

Recent research has explored methods for extracting personal attributes
from unstructured text. A significant portion of this work has been
specifically focused on dialogues. Wang et al.~(2022) explored the use
of a language model to extract and infer personal attributes from
human-human dialogue. DeLucia et al.~(2024) presented a method for
adapting persona extraction models to diverse conversational settings,
enhancing the quality of persona information extracted from dialogues
and reducing the need for extensive human annotation. Wu et al.~(2024)
introduced In-Dialogue Learning, a framework that enhances pre-trained
language models to utilize dialogue history for persona extraction in
personalized dialogue generation

All aforementioned studies employ triplet format. This means the format
(s, r, o), where the relation (r) indicates the persona attribute type
of the subject (s) towards the object (o). The benefits of triplet
format are 1, 2, and 3.

Despite their promising results, these studies have limitations. They
rely on predefined lists of attributes, which constrain their ability to
capture a broad range of relevant information and limit flexibility. In
addition, models need to be fine-tuned to address these constraints,
especially when applied to product reviews. There is a need for data to
finetune.

\hypertarget{approach}{%
\section{Approach}\label{approach}}

In line with previous research on extracting personal attributes, this
study will also use the triplet format to store these attributes.
However, our approach diverges by addressing some of the key limitations
of earlier studies, including that previous research often relied on
predefined lists of attributes. Inspired by recent advancements in
pre-trained large language models (LLMs), this work introduces a novel
approach for extracting personal attributes using these models. Our
method aims to extract triplets from product reviews without being
limited by predefined attribute categories, thus offering a more
comprehensive and adaptable solution to capture the various personal
attributes revealed in customer feedback.

We present an end-to-end pipeline, starting from raw text all the way
through to a constructed knowledge graph. The approach , text
pre-processing, triplet extraction post-processing.

\hypertarget{main-contributions}{%
\section{Main Contributions}\label{main-contributions}}

The main contributions of this thesis are as follows: * Pipeline: We
developed an end-to-end pipeline that leverages large language models
(LLMs) to automate the construction of Knowledge Graphs from
unstructured data. * Comparison: We provided a comparison between our
LLM-based approach and traditional triplet extraction methods. * Case
Study Implementation: We conducted a detailed case study to evaluate the
performance of our framework. * Future Direction: We identified
potential areas for future research and development, offering insights
into how the end-to-end pipeline can be further improved and extended.

\hypertarget{thesis-outline}{%
\section{Thesis Outline}\label{thesis-outline}}

Placeholder for thesis outline

\bookmarksetup{startatroot}

\hypertarget{system-architecture}{%
\chapter{System architecture}\label{system-architecture}}

Several different software components are required to ensure the system
can:

\begin{itemize}
\tightlist
\item
  persist general data on disk, such as node telemetry and hardware
  information
\item
  receive and store large volumes of radio signal samples organized in
  \enquote{campaigns}
\item
  maintain a reliable and fault-tolerant connection between the backend
  and an unknown number of nodes
\end{itemize}

The currently implemented software architecture is summarized in
Figure~\ref{fig-systemArch}.

\begin{figure}[h!]

{\centering \includegraphics[width=\textwidth,height=0.4\textheight]{index_files/mediabag/28a0707630a26b7937bb55ffeebfc4d2a7a67a2a.pdf}

}

\caption{\label{fig-systemArch}high-level system architecture}

\end{figure}

\hypertarget{high-level-components}{%
\section{High-level components}\label{high-level-components}}

OpenRFSense is built upon several different modern technologies,
interconnected by mainstream networking protocols.

\hypertarget{hardware}{%
\subsection{Hardware}\label{hardware}}

OpenRFSense is designed to run on a variety of hardware platforms,
including PCs, servers, and embedded systems. The hardware requirements
depend on the specific use case and workload, but typically include a
multi-core processor, sufficient memory and storage, and one or more
software-defined radio (SDR) devices acting as signal receivers.

\hypertarget{software-defined-radio}{%
\subsection{Software-defined radio}\label{software-defined-radio}}

OpenRFSense relies on one or more SDR devices to capture and analyze
radio frequency signals. The SDR devices are connected to the hardware
platform via USB or PCIe interfaces and can be controlled using
open-source software libraries such as GNU Radio. The SDR device is
responsible for capturing radio signals from the environment.

\hypertarget{operating-system}{%
\subsection{Operating system}\label{operating-system}}

Thanks to the choice of programming language (Golang
\protect\hyperlink{ref-GoProgrammingLanguage}{{[}1{]}}), both the
backend and node software can be compiled and deployed on most modern
operating systems and hardware configurations (e.g.~Windows on ARM,
Linux on x86, MacOS on M1). The backend service is also developed with
containerization in mind, making it possible to deploy the software
securely and avoiding strict dependency on host system configuration.

\hypertarget{signal-processing-and-analysis}{%
\subsection{Signal processing and
analysis}\label{signal-processing-and-analysis}}

OpenRFSense uses a variety of signal processing and analysis techniques
to extract useful information from the captured radio frequency signals.
This includes techniques such as Fast Fourier Transform
\protect\hyperlink{ref-heidemanGaussHistoryFast1984}{{[}2{]}} (FFT),
digital signal processing (DSP), and machine learning algorithms. See
Section~\ref{sec-signalProcessing} for a more in-depth explanation of
the data analysis process.

\hypertarget{user-interface-and-management}{%
\subsection{User interface and
management}\label{user-interface-and-management}}

OpenRFSense provides a web-based user interface for configuring,
managing, and visualizing the captured data. The user interface includes
features such as real-time spectrum displays, signal analysis tools, and
alerting mechanisms. The system also provides APIs for integration with
external tools and services.

\hypertarget{backend-architecture}{%
\section{Backend architecture}\label{backend-architecture}}

The backend service is a monolithic application which combines user
interface and data management code. It can be broken down into several
modular components which all depend only on shared configuration in the
form of a file or environment variables on the host system.

The internal messaging service maintains a constant communication stream
between the backend and the various nodes deployed by the user. It
provides a scalable and fault-tolerant solution for message delivery to
the remote nodes. It is documented in
Chapter~\ref{sec-serviceMessaging}.

Another critical component is the web-based user interface (UI) which is
the main interactive mean for the user to communicate with the other
components inside the backend service. An in-depth explanation, complete
with other high-level considerations, is contained in
Chapter~\ref{sec-ui}.

\hypertarget{sec-restApi}{%
\subsection{REST API}\label{sec-restApi}}

To allow external access to stored data, a REST (REpresentational State
Transfer, defined in \protect\hyperlink{ref-fieldingArch2000}{{[}3{]}})
API is provided by the backend. Such a system allows authorized
applications to query data through a standardized interface. Currently,
the following data can be fetched from the backend:

\begin{itemize}
\tightlist
\item
  signal measurements taken by a specific node and belonging to a
  certain campaign
\item
  a list of all nodes currently connected to the messaging system
\item
  node metrics and system status for a specific node
\end{itemize}

The API can also actively request actions to be carried out by the
nodes. To ensure bad actors cannot arbitrarily send command requests to
the system, all requests which perform an action that can change the
state of the system require a form of authentication. The currently
implemented authentication method is Basic HTTP Authentication (a Web
standard, see \protect\hyperlink{ref-RFC7235}{{[}4{]}}).

\hypertarget{sec-signalStorage}{%
\subsection{Signal measurements storage}\label{sec-signalStorage}}

Since radio signal data can be dense and structured as large inbound TCP
packets, a storage layer capable of extremely fast writes to memory is
needed. The BadgerDB \protect\hyperlink{ref-BadgerDB}{{[}5{]}} key-value
store was chosen due to the maturity of the software and it being
natively developed in Golang, providing easier integration into the
existing backend code. BadgerDB is capable, in optimal conditions
(sufficient RAM and a modern solid state disk), of writing several
gigabytes of data per second to disk. In practical testing, BadgerDB has
proven more than sufficient for handling raw inbound signal data from
several nodes at a time.

\hypertarget{database-access-layer}{%
\subsection{Database access layer}\label{database-access-layer}}

A support database is used to store persistent data, which is required
to survive service outages or is generally better kept for an
unspecified amount of time. PostgreSQL
\protect\hyperlink{ref-groupPostgreSQL2023}{{[}6{]}} was chosen mainly
due to the maturity and notoriety of the project. From the project
website:

\begin{quote}
\emph{PostgreSQL is a powerful, open source object-relational database
system that uses and extends the SQL language combined with many
features that safely store and scale the most complicated data
workloads.}
\end{quote}

PostgreSQL is widely available and easily deployable even using
container technology. A native Golang implementation of the database
connection and communication protocol is used internally to store data
such as past node metrics received through the messaging system, or
signal measurement campaign data (start date and time, end date and
time, etc.). Such data is then available for querying through the REST
API documented above (in Section~\ref{sec-restApi}).

\hypertarget{node-architecture}{%
\section{Node architecture}\label{node-architecture}}

The node software is developed upon much the same technologies as the
backend to ease integration between the two and reduce development
overhead. Being developed to run natively on embedded devices as a
self-contained application, several software components are needed to
fullfil all required functions.

The application is available both as a standalone, statically-linked
executable and a more user-friendly prebuilt system image based on the
Raspbian OS. The system image comes pre-configured with required system
components such as:

\begin{itemize}
\tightlist
\item
  a SystemD \protect\hyperlink{ref-Systemd}{{[}7{]}} service which
  starts the node management software on boot
\item
  the OpenRFSense node management application
\item
  the external, low-level program which interfaces with the SDR hardware
  (\texttt{orfs-sensor})
\item
  some useful signal decoding software, also part of the project
\end{itemize}

The following sections describe the main components of the node
management software.

\hypertarget{system-metrics-collector}{%
\subsection{System metrics collector}\label{system-metrics-collector}}

A critical component for the node software is the metrics collector.
This component is responsible for providing a current snapshot of the
status of both hardware and software of the system, which will then be
stored by the backend. The metrics message format and message exchange
is documented more extensively in Section~\ref{sec-metrics}.

\hypertarget{process-manager}{%
\subsection{Process manager}\label{process-manager}}

The node software relies on an external process for data acquisition and
delivery. Not being originally developed with integration in mind, such
a program is managed as a zero-knowledge black box, executed as a child
process by generating command-line arguments and passing them to the
program.

\begin{figure}[h!]

{\centering \includegraphics[width=0.7\textwidth,height=\textheight]{index_files/mediabag/2ede55a8c9c0223312f5fe980e0c838aa78c9b6b.pdf}

}

\caption{measurement process flow}

\end{figure}

\hypertarget{user-interface}{%
\subsection{User interface}\label{user-interface}}

Much like the backend, the node provides users with a web-based UI to
manage system functionality. Additionally, the UI lets users configure
and reboot the system without needing an external monitor or peripherals
for easier system administration of embedded devices. A description of
the node UI with screenshots is contained in Section~\ref{sec-nodeUi}.

The interface is only accessible on the same network as the node, to
avoid exposing administration functionality to the public internet.
Additionally, for thirty minutes after the system is booted, the UI is
accessible by connecting to the device through a temporary WiFi access
point if the hardware is capable of creating one such connection.

\bookmarksetup{startatroot}

\hypertarget{sec-serviceMessaging}{%
\chapter{Service messaging}\label{sec-serviceMessaging}}

In OpenRFSense, NATS messaging \protect\hyperlink{ref-NATS}{{[}8{]}} is
used to enable communication between the backend and the nodes. NATS
(Neural Autonomic Transport System) is a lightweight messaging system
that is commonly used for machine-to-machine (M2M) communication in IoT
(Internet of Things) applications. A messaging system such as NATS is
based on message delivery guarantees and optional message retention for
later delivery. These capabilities ensure service message resiliency
through high-latency connections and in-flight data loss.

The backend acts as the NATS server and the nodes act as NATS clients
which connect to the server. When a node is powered on, it connects to
the NATS server and subscribes to a \enquote{subject}, which is a named
destination to which messages are published. The node can subscribe to
one or more subjects depending on its capabilities. Some subjects
currently being used for service messages are:

\begin{itemize}
\tightlist
\item
  \texttt{node.all} for node system metrics (see
  Section~\ref{sec-metrics})
\item
  \texttt{node.all.aggregated} and \texttt{node.all.raw} for
  backend-to-node signal recording commands (see
  Section~\ref{sec-commands})
\item
  \texttt{node.\$hardware\_id.\$command}, where \texttt{\$hardware\_id}
  is the node's unique identifier and \texttt{\$command} is a specific
  action being requested
\end{itemize}

\hypertarget{sec-metrics}{%
\section{Node metrics}\label{sec-metrics}}

The backend, when requested through the REST API or the web UI, sends a
\enquote{status request} message to a specific NATS subject (to which
all nodes are required to listen on). All nodes then start collecting
internal metrics (such as CPU usage and temperature, memory usage etc.)
and respond on the same subject with the collected data, encoded in JSON
(JavaScript Object Notation) format
\protect\hyperlink{ref-ecmaECMA404JSONData2017}{{[}9{]}}.

\begin{figure}[h!]

{\centering \includegraphics{index_files/mediabag/c498518269b754f0cb52da425ed2e89a40bb3797.pdf}

}

\caption{\label{fig-metrics}node metrics request flow}

\end{figure}

\hypertarget{metrics-contents}{%
\subsection{Metrics contents}\label{metrics-contents}}

A metrics object contains useful system data such as:

\begin{itemize}
\tightlist
\item
  the system hostname
\item
  a unique hardware ID derived from the motherboard model
\item
  the motherboard/hardware model of the system
\item
  the system's online time (uptime) in milliseconds since it was powered
  on
\end{itemize}

Additionally, a variable-content section is reserved for less crucial
information. This section is never required nor checked by the backend,
but its data, if present, is shown to the user through the web UI.
Currently, such data is provided by small standardized software modules
(\enquote{providers}), which are called by the node management software
if present. Some currently implemented providers are:

\begin{itemize}
\tightlist
\item
  node geographical location in GeoJSON
  \protect\hyperlink{ref-RFC7946}{{[}10{]}} format
\item
  filesystem usage
\item
  memory (RAM) usage
\item
  network connections
\item
  current node availability (free, busy, handling errors, etc.)
\end{itemize}

\hypertarget{message-encoding}{%
\subsection{Message encoding}\label{message-encoding}}

As stated above, node metrics are encoded in JSON format to be easily
accessible by external services, since JSON is a well-established
industry standard data format. An example of the metrics data encoded in
textual JSON is Listing~\ref{lst-metrics}.

\begin{codelisting}

\caption{example metrics object in JSON format}

\hypertarget{lst-metrics}{%
\label{lst-metrics}}%
\begin{Shaded}
\begin{Highlighting}[numbers=left,,]
\FunctionTok{\{}
  \DataTypeTok{"id"}\FunctionTok{:} \StringTok{"kgslnximugwhsfnwbjknwbv"}\FunctionTok{,}
  \DataTypeTok{"hostname"}\FunctionTok{:} \StringTok{"raspberry"}\FunctionTok{,}
  \DataTypeTok{"model"}\FunctionTok{:} \StringTok{"Raspberry Pi 3B"}\FunctionTok{,}
  \DataTypeTok{"uptime"}\FunctionTok{:} \DecValTok{186339250000000}\FunctionTok{,}
  \DataTypeTok{"providers"}\FunctionTok{:} \OtherTok{[]}
\FunctionTok{\}}
\end{Highlighting}
\end{Shaded}

\end{codelisting}

\hypertarget{sec-commands}{%
\section{Commands}\label{sec-commands}}

The backend can also publish messages to node-specific subjects. This
enables the backend to remotely control the nodes and send configuration
updates or commands. The messages can be one of two types:

\begin{enumerate}
\def\labelenumi{\arabic{enumi}.}
\tightlist
\item
  empty messages sent on a specific NATS subject (\enquote{pings})
\item
  JSON-encoded messages containing relevant context for the requested
  action
\end{enumerate}

When a node receives a remote command message, it processes the command
and sends a response message back to the backend. The response message
can contain any relevant information about the status of the command,
such as success or failure, and any associated data. Generally, nodes
send their system metrics as response but it is not a strict
requirement.

\begin{figure}[h!]

{\centering \includegraphics{index_files/mediabag/594ec6cc8bb7a28e51cf7a27e0d8aa69c49e0602.pdf}

}

\caption{\label{fig-commands}sensor campaign command flow}

\end{figure}

\bookmarksetup{startatroot}

\hypertarget{sec-ui}{%
\chapter{User interface}\label{sec-ui}}

The user interface plays an essential role in enabling users to interact
with the system effectively. A well-designed user interface can make it
easy for users to navigate between pages and perform several complex
tasks related to node system and data management.

The UI is served to the user as a collection of web pages by the backend
service. The pages are generated on demand with up-to-date information
and take advantage of modern web technologies.

\hypertarget{design-principles}{%
\section{Design principles}\label{design-principles}}

Several design principles guided the development of the user interface,
mainly:

\begin{enumerate}
\def\labelenumi{\arabic{enumi}.}
\tightlist
\item
  Consistency: Elements such as fonts, colors, and icons are consistent
  throughout the interface, making it easier for users to navigate
  between UI sections and use the system.
\item
  Simplicity: The user interface is simple and easy to use. The
  simplicity principle ensures that the interface is not cluttered with
  unnecessary elements, making it easy for users to find what they need.
\item
  Feedback: The user interface provides users with immediate feedback to
  let them know that their actions have been registered. This principle
  ensures that users are not left wondering whether their actions have
  been successful.
\item
  Accessibility: The user interface is designed to be accessible to all
  users, including those with disabilities. This principle ensures that
  users with different abilities can use the system without encountering
  any usability barriers.
\end{enumerate}

\newpage{}

\hypertarget{interface-components-and-structure}{%
\section{Interface components and
structure}\label{interface-components-and-structure}}

The UI is aimed at users with a certain technical background and
previous knowledge of the system, but it is designed to be easily
navigated through by any given user.

\begin{figure}[h!]

{\centering \includegraphics{index_files/mediabag/05e1db55b179033dd01d6bbbcb09b1917c1dda84.pdf}

}

\caption{UI hyperlink structure}

\end{figure}

\hypertarget{dashboard}{%
\subsection{Dashboard}\label{dashboard}}

The dashboard acts as the homepage of the UI. It is the first page
served by the backend on the root address and is the first interactive
element the user will be shown on connection to the system.

\begin{figure}[h!]

{\centering \includegraphics[width=0.8\textwidth,height=\textheight]{assets/screenshots/dashboard.pdf}

}

\caption{Dashboard}

\end{figure}

The first element presented by the UI is a table containing useful
information about all nodes connected to the backend, such as node
location and unique identifier. Each node can be selected with a
checkbox to take part in a measurement campaign (configurable through a
pop-up modal dialog). Every row also contains a link to a page which
visualizes a node's metrics and location in more detail (see
Section~\ref{sec-nodeOverview}).

The homepage of the UI provides a map of the world, with the locations
of all the sensors in the OpenRFSense network indicated by markers.
Clicking on the markers yields basic information about each node, such
as its ID and the last time it was active, and a link the node-specific
overview page.

\hypertarget{sec-nodeOverview}{%
\subsection{Node overview page}\label{sec-nodeOverview}}

A node-specific overview page is also generated for each node connected
to the backend and the messaging system. This page generally contains
all required node metrics returned by a request to the node, as well as
more detailed node system information if available (see
Section~\ref{sec-metrics}). The node overview page also contains a map
but only showing a single marker corresponding to the node being queried
and centered on the node's location, if provided.

\begin{figure}[h!]

{\centering \includegraphics[width=0.8\textwidth,height=\textheight]{assets/screenshots/node.pdf}

}

\caption{Node overview page}

\end{figure}

A table containing a list of all past and ongoing measurement campaigns
is also provided to the user, with management controls on each row to
download the signal data in various formats or delete it from the
storage. The controls are non-interactive if the campaign is still
ongoing.

\newpage{}

\hypertarget{api-documentation}{%
\subsection{API documentation}\label{api-documentation}}

To make the REST API (defined in Section~\ref{sec-restApi}) easily
accessible to users and external applications, a documentation page is
automatically generated using Swagger (OpenAPI specification version 2.0
\protect\hyperlink{ref-OpenAPISpecification2}{{[}11{]}}). The generated
page contains all the necessary documentation to query data and send
commands to the backend, complete with examples and request and response
formats in textual JSON.

\begin{figure}[h!]

{\centering \includegraphics[width=0.8\textwidth,height=\textheight]{assets/screenshots/api-docs.pdf}

}

\caption{API documentation}

\end{figure}

\newpage{}

\hypertarget{sec-nodeUi}{%
\subsection{Node system interface}\label{sec-nodeUi}}

Each node's local management UI provides a high-level management over
the host system running the node. As such, it must present important
information in logical order, starting with connection status (both WiFi
and physical). On remote devices, this facilitates troubleshooting and
provides an alternative method of restoring connection to a node.

\begin{figure}[h!]

{\centering \includegraphics[width=0.8\textwidth,height=\textheight]{assets/screenshots/node-ui.png}

}

\caption{Node system management page}

\end{figure}

A simple text editor is also provided for the node management software
configuration. Each change followed by a click on the \enquote{Save}
button will write the changes to the configuration file and reboot the
system.

\newpage{}

\bookmarksetup{startatroot}

\hypertarget{data-collection-and-analysis}{%
\chapter{Data collection and
analysis}\label{data-collection-and-analysis}}

Data collection is at the core of the OpenRFSense project and, as such,
requires some complex and specialized components. Overall, the data
collection procedure can be better explained as an ordered pipeline,
which starts from the remote node.

\hypertarget{data-acquisition}{%
\section{Data acquisition}\label{data-acquisition}}

The node software utilizes a set of low-level hardware APIs to interface
with the SDR receiver(s) and collect the RF signals. The main data
collection routine is handled by a customized version of a pre-existing
Electrosense software, which interfaces with connected SDR devices and
streams data towards an external destination.

The data is collected in one of two forms:

\begin{itemize}
\tightlist
\item
  raw spectral density ordered by frequency
\item
  spectral density averaged over an arbitrary number of seconds using
  Fast Fourier Transform (FFT)
\end{itemize}

\hypertarget{sec-signalProcessing}{%
\section{Signal processing}\label{sec-signalProcessing}}

The signal processing component is responsible for cleaning up the
collected RF signals and extracting the relevant features for analysis.
The component utilizes a set of algorithms to filter out noise and
interference from the collected signals. Once the signals are cleaned
up, the component extracts relevant features such as the signal
strength, modulation type, and frequency. At the time of writing, the
following signal encodings can be preprocessed:

\begin{itemize}
\tightlist
\item
  Aircraft Communications Addressing and Reporting System (ACARS)
\item
  LTE, both User Equipment-to-cell and cell-to-User Equipment
\item
  Mode S
  \protect\hyperlink{ref-internationalcivilaviationorganizationManualModeSpecific2004}{{[}12{]}}
  collision avoidance messages
\item
  Aeronautical Information Service (AIS) air traffic information
\item
  AM/FM radio
\end{itemize}

\hypertarget{data-streaming}{%
\section{Data streaming}\label{data-streaming}}

The processed data is serialized in a JSON-based binary format (Apache
Avro \protect\hyperlink{ref-ApacheAvroDocumentation}{{[}13{]}}),
together with extra metadata such as:

\begin{itemize}
\tightlist
\item
  time of recording
\item
  node hardware identifier and campaign identifier code
\item
  sensor hardware configuration (center frequency, antenna gain,
  estimated noise floor, etc.)
\end{itemize}

The encoded packets are then sent over the Internet to the backend
service through several raw TCP connections per receiver to maximize
flow. Surprisingly, even low-cost modern SDR receivers tend to output
large quantities of data (usually in the order of several megabytes per
second). This requires the usage of clever throughput maximization
algorithms such as TCP Fast Open
\protect\hyperlink{ref-radhakrishnanTCPFastOpen2011}{{[}14{]}} and
\texttt{SO\_REUSEPORT}.

\hypertarget{data-storage}{%
\section{Data storage}\label{data-storage}}

The cleaned up and feature-extracted data is received by the backend and
partially deserialized to extract the node hardware identifier and
campaign identifier code, which are used to derive a unique key for that
packet to be used in storage. The binary-encoded data is then stored in
a database embedded in the backend for later analysis. The data is
written to disk in a structured format that allows for fast querying and
analysis. The storage software component is explained in more detail in
Section~\ref{sec-signalStorage}.

\hypertarget{analysis-of-spectrum-data}{%
\section{Analysis of spectrum data}\label{analysis-of-spectrum-data}}

To analyze the collected spectrum data, several techniques can be
employed. One approach is to use statistical methods to identify
patterns or anomalies in the data. For example, signal strength and
frequency distribution can be analyzed to detect unusual patterns, which
may indicate the presence of interference sources. Another approach is
to use machine learning algorithms to automatically classify different
types of signals based on their frequency, bandwidth, and modulation.
This can help to identify different types of signals, such as WiFi,
Bluetooth, or cellular signals, and to detect any new or unknown
signals.

Overall, the collected spectrum data is a valuable resource for a wide
range of applications. The different techniques for analyzing the data
can provide useful insights into the electromagnetic environment, giving
users easy access to dense and reasonably accurate research samples.

\bookmarksetup{startatroot}

\hypertarget{conclusion-and-future-work}{%
\chapter{Conclusion and future work}\label{conclusion-and-future-work}}

OpenRFSense was developed as an experimental software suite powered by
modern software-defined radio technology. The following sections contain
some closing remarks and reflections about the project and the status of
the SDR-based data analysis as a whole.

\hypertarget{evaluation-of-system-performance-and-limitations}{%
\section{Evaluation of system performance and
limitations}\label{evaluation-of-system-performance-and-limitations}}

As with any project, OpenRFSense has its limitations that must be
acknowledged. Below are some of the limitations of the project.

\begin{enumerate}
\def\labelenumi{\arabic{enumi}.}
\tightlist
\item
  Coverage: OpenRFSense relies on users who can install and maintain
  their own sensors. As a result, there may be regions where the sensor
  density is low, leading to gaps in the collected data.
\item
  Sensor accuracy: The accuracy of the collected data is dependent on
  the quality of the sensors installed. Low-quality sensors may produce
  unreliable data that could affect the accuracy of the analysis and
  interpretation of the collected data.
\item
  Signal interference: Interference from other radio sources can tamper
  with the accuracy of the collected data. For example, electromagnetic
  radiation from other devices or sources could create a noisy signal
  that can be challenging to analyze.
\item
  Cost: The cost of installing and maintaining the nodes could be a
  restraining factor, especially in regions where resources are limited.
\item
  Security: The sensors collect data from the environment, and ensuring
  the security and privacy of the collected data is crucial. The project
  must ensure that the data collected is not misused or accessed by
  unauthorized individuals.
\item
  Data quality: The data collected may not always be of high quality due
  to a variety of reasons such as environmental conditions, sensor
  errors, and data transmission errors. Therefore, it is essential to
  have mechanisms in place to ensure data quality control and filtering.
\end{enumerate}

\hypertarget{applications-of-the-project}{%
\section{Applications of the
project}\label{applications-of-the-project}}

The collected spectrum data can be used for a wide range of
applications, including radio frequency interference (RFI)
identification, spectrum occupancy measurements, and wireless network
planning.

RFI identification is a major application of spectrum data. OpenRFSense
provides a real-time view of the electromagnetic environment, allowing
users to detect and locate interference sources. The data can also be
used for spectrum monitoring, such as measuring the spectrum occupancy
of various frequencies. This information can be used to identify
underutilized spectrum bands, which can then be allocated for new
services or applications.

Another important application of the data is wireless network planning.
OpenRFSense provides accurate and up-to-date information on the
electromagnetic environment, allowing users to identify optimal
locations for wireless access points or base stations. The data can also
be used to optimize the configuration of said networks, such as
adjusting the transmission power or selecting the most appropriate
frequency band.

Hobbyists and researchers have also been able to carry out more complex
tasks thanks to SDR receivers, such as satellite tracking and image
reception
\protect\hyperlink{ref-peraltaSatelliteTelemetryImage2018}{{[}15{]}}.
This proves SDR-sourced data to be useful even for longer range and
higher gain spectrum analysis.

\hypertarget{future-directions-for-research-and-development}{%
\section{Future directions for research and
development}\label{future-directions-for-research-and-development}}

Being released as open-source software, OpenRFSense is not intended to
be a finished, industrial-grade product. Some future areas in which the
project could be improved or may need additional research are the
following.

\begin{enumerate}
\def\labelenumi{\arabic{enumi}.}
\tightlist
\item
  Integration with other sensor networks: OpenRFSense can be integrated
  with other sensor networks, such as those monitoring air quality,
  water quality, and weather. Integration with these networks can
  provide a more comprehensive picture of the environment and the impact
  of electromagnetic radiation on it.
\item
  Machine learning and AI: OpenRFSense can be enhanced with machine
  learning and artificial intelligence (AI) algorithms to improve the
  accuracy and speed of data analysis. These algorithms can help
  identify patterns and anomalies in the data and assist in making
  predictions and recommendations based on the data.
\item
  Real-time monitoring: The current version of OpenRFSense only provides
  playback of recorded spectrum data. However, further development can
  be done to improve the latency and accuracy of the system, at the
  point of providing users with real-time signal visualization. This can
  require the use of more advanced sensors and improved data
  transmission technologies.
\item
  Visualization and data analysis tools: OpenRFSense generates a large
  amount of data, and tools are needed to help users visualize and
  analyze the data. This can involve the development of new
  visualization techniques, such as heat maps and trend analysis, as
  well as ad-hoc data mining tools and other task-specific analysis
  tools.
\item
  Standardization and data sharing: OpenRFSense can benefit from
  standardization of data collection and sharing protocols, such as a
  universal format for spectrum data storage. This can help ensure that
  data collected by different sensors and networks can be easily shared
  and combined to provide a more comprehensive picture of the
  environment.
\end{enumerate}

\appendix

\titleformat{\chapter}{\normalfont\Huge\bfseries}{Appendix \thechapter}{1em}{}

\bookmarksetup{startatroot}

\hypertarget{code-structure}{%
\chapter{Code structure}\label{code-structure}}

All the code for the OpenRFSense project is hosted on Github, a free Git
repository hosting service with additional code and community management
tools. At the time of writing, the various components are placed in the
following repositories under the
\href{https://github.com/openrfsense}{OpenRFSense organization}:

\begin{itemize}
\tightlist
\item
  \texttt{backend}: source code and tooling for the backend service
\item
  \texttt{node}: source code and tooling for the node management
  software
\item
  \texttt{orfs-sensor}: source code for the signal recording software
\item
  \texttt{image}: various scripts and tooling to generate the
  pre-configured system image for Raspberry Pi boards
\item
  \texttt{openrfsense.github.io}: technical documentation for the
  project (see Appendix~\ref{sec-techDoc})
\item
  \texttt{common}: shared code which would otherwise be duplicated in
  both node and backend
\end{itemize}

\bookmarksetup{startatroot}

\hypertarget{project-branding}{%
\chapter{Project branding}\label{project-branding}}

Special care was put behind giving an identity to the OpenRFSense
project through graphic design. A simple but meaningful logo contains
the abbreviated project name and a round symbology reminiscent of
WiFi-associated commercial products.

\begin{figure}[h!]

{\centering \includegraphics[width=0.6\textwidth,height=\textheight]{index_files/mediabag/assets/logos/openrf.pdf}

}

\caption{Official OpenRFSense logo}

\end{figure}

The color palette has enough contrast to be accessible but maintains at
least one saturated color.

\begin{figure}[h!]

{\centering \includegraphics{index_files/mediabag/assets/palette.pdf}

}

\caption{Color palette}

\end{figure}

\bookmarksetup{startatroot}

\hypertarget{sec-techDoc}{%
\chapter{Technical documentation}\label{sec-techDoc}}

All software components for the OpenRFSense project contain
self-documenting code with numerous comments and additional
language-specific documentation tooling (\texttt{godoc}). Moreover,
additional user-oriented technical documentation has been written and is
deployed automatically on Github's static file hosting service (Github
Pages).

The documentation is written in Markdown format and is rendered to a
static website using the Jekyll program. It includes code blocks with
ready-to-use examples to prepare deployments and configure the services
before starting the system, together with overviews of the various
components. All pages also feature a search widget which facilitates
navigation through the text contents.

\begin{figure}[h!]

{\centering \includegraphics[width=0.8\textwidth,height=\textheight]{assets/screenshots/docs.pdf}

}

\caption{OpenRFSense documentation}

\end{figure}

\bookmarksetup{startatroot}

\hypertarget{bibliography}{%
\chapter*{Bibliography}\label{bibliography}}
\addcontentsline{toc}{chapter}{Bibliography}

\markboth{Bibliography}{Bibliography}

\hypertarget{refs}{}
\begin{CSLReferences}{0}{0}
\leavevmode\vadjust pre{\hypertarget{ref-GoProgrammingLanguage}{}}%
\CSLLeftMargin{{[}1{]} }%
\CSLRightInline{{``The {Go Programming Language}.''} {[}Online{]}.
Available: \url{https://go.dev/}}

\leavevmode\vadjust pre{\hypertarget{ref-heidemanGaussHistoryFast1984}{}}%
\CSLLeftMargin{{[}2{]} }%
\CSLRightInline{M. Heideman, D. Johnson, and C. Burrus, {``Gauss and the
history of the fast fourier transform,''} \emph{IEEE ASSP Mag.}, vol. 1,
no. 4, pp. 14--21, Oct. 1984, doi:
\href{https://doi.org/10.1109/MASSP.1984.1162257}{10.1109/MASSP.1984.1162257}.
{[}Online{]}. Available:
\url{http://ieeexplore.ieee.org/document/1162257/}. {[}Accessed: Feb.
27, 2023{]}}

\leavevmode\vadjust pre{\hypertarget{ref-fieldingArch2000}{}}%
\CSLLeftMargin{{[}3{]} }%
\CSLRightInline{R. T. Fielding and R. N. Taylor, {``Architectural styles
and the design of network-based software architectures,''} {University
of California, Irvine}, 2000. }

\leavevmode\vadjust pre{\hypertarget{ref-RFC7235}{}}%
\CSLLeftMargin{{[}4{]} }%
\CSLRightInline{R. Fielding and J. Reschke, {``Hypertext transfer
protocol ({HTTP}/1.1): {Authentication},''} {RFC Editor / RFC Editor};
Internet Requests for Comments, RFC 7235, Jun. 2014 {[}Online{]}.
Available: \url{http://www.rfc-editor.org/rfc/rfc7235.txt}}

\leavevmode\vadjust pre{\hypertarget{ref-BadgerDB}{}}%
\CSLLeftMargin{{[}5{]} }%
\CSLRightInline{{``{BadgerDB}.''} {[}Online{]}. Available:
\url{https://dgraph.io/docs/badger}}

\leavevmode\vadjust pre{\hypertarget{ref-groupPostgreSQL2023}{}}%
\CSLLeftMargin{{[}6{]} }%
\CSLRightInline{P. G. D. Group, {``{PostgreSQL},''} Feb. 25, 2023.
{[}Online{]}. Available: \url{https://www.postgresql.org/}}

\leavevmode\vadjust pre{\hypertarget{ref-Systemd}{}}%
\CSLLeftMargin{{[}7{]} }%
\CSLRightInline{{``Systemd.''} {[}Online{]}. Available:
\url{https://www.freedesktop.org/wiki/Software/systemd/}}

\leavevmode\vadjust pre{\hypertarget{ref-NATS}{}}%
\CSLLeftMargin{{[}8{]} }%
\CSLRightInline{{``{NATS}.''} {[}Online{]}. Available:
\url{https://nats.io/about}}

\leavevmode\vadjust pre{\hypertarget{ref-ecmaECMA404JSONData2017}{}}%
\CSLLeftMargin{{[}9{]} }%
\CSLRightInline{ECMA, \emph{{ECMA-404}: {The JSON} data interchange
syntax}. {Geneva, Switzerland}: {ECMA (European Association for
Standardizing Information and Communication Systems)}, 2017
{[}Online{]}. Available:
\url{https://www.ecma-international.org/publications-and-standards/standards/ecma-404/}}

\leavevmode\vadjust pre{\hypertarget{ref-RFC7946}{}}%
\CSLLeftMargin{{[}10{]} }%
\CSLRightInline{H. Butler, M. Daly, A. Doyle, S. Gillies, S. Hagen, and
T. Schaub, {``The {GeoJSON} format,''} {RFC Editor / RFC Editor};
Internet Requests for Comments, RFC 7946, Aug. 2016. }

\leavevmode\vadjust pre{\hypertarget{ref-OpenAPISpecification2}{}}%
\CSLLeftMargin{{[}11{]} }%
\CSLRightInline{{``{OpenAPI Specification} - {Version} 2.0 \textbar{}
{Swagger}.''} {[}Online{]}. Available:
\url{https://swagger.io/specification/v2/}}

\leavevmode\vadjust pre{\hypertarget{ref-internationalcivilaviationorganizationManualModeSpecific2004}{}}%
\CSLLeftMargin{{[}12{]} }%
\CSLRightInline{I. C. A. Organization, Ed., \emph{Manual on mode {S}
specific services}, 2nd. ed. {Montreal}: {ICAO}, 2004. }

\leavevmode\vadjust pre{\hypertarget{ref-ApacheAvroDocumentation}{}}%
\CSLLeftMargin{{[}13{]} }%
\CSLRightInline{{``Apache {Avro Documentation}.''} {[}Online{]}.
Available: \url{https://avro.apache.org/docs/}}

\leavevmode\vadjust pre{\hypertarget{ref-radhakrishnanTCPFastOpen2011}{}}%
\CSLLeftMargin{{[}14{]} }%
\CSLRightInline{S. Radhakrishnan, Y. Cheng, J. Chu, A. Jain, and B.
Raghavan, {``{TCP} fast open,''} in \emph{Proceedings of the {Seventh
COnference} on emerging {Networking EXperiments} and {Technologies}},
Dec. 2011, pp. 1--12, doi:
\href{https://doi.org/10.1145/2079296.2079317}{10.1145/2079296.2079317}
{[}Online{]}. Available:
\url{https://dl.acm.org/doi/10.1145/2079296.2079317}}

\leavevmode\vadjust pre{\hypertarget{ref-peraltaSatelliteTelemetryImage2018}{}}%
\CSLLeftMargin{{[}15{]} }%
\CSLRightInline{D. J. M. Peralta, D. S. D. Santos, A. Tikami, W. A. D.
Santos, and E. W. R. Pereira, {``Satellite {Telemetry} and {Image
Reception} with {Software Defined Radio Applied} to {Space Outreach
Projects} in {Brazil},''} \emph{An. Acad. Bras. Ciênc.}, vol. 90, no. 3,
pp. 3175--3184, Sep. 2018, doi:
\href{https://doi.org/10.1590/0001-3765201820170955}{10.1590/0001-3765201820170955}.
{[}Online{]}. Available:
\url{http://www.scielo.br/scielo.php?script=sci_arttext\&pid=S0001-37652018000603175\&tlng=en}}

\end{CSLReferences}
\endgroup
\clearpage


\backmatter

\clearpage

\end{document}
